% Copyright (C) 2026, Lux Industries Inc. All rights reserved.
% Lux DEX cross-chain HTLC atomic-swap settlement — construction, precompile
% interface, and threat model. Replaces the custodial lock-and-mint bridge and
% the dexvm custodial proxy. This document is normative: implement from it.
\documentclass[11pt]{article}
\usepackage[margin=1in]{geometry}
\usepackage{amsmath,amssymb,amsthm}
\usepackage{enumitem}
\usepackage{booktabs}
\usepackage{listings}
\usepackage{hyperref}
\usepackage{xcolor}

\lstset{basicstyle=\ttfamily\small,breaklines=true,columns=fullflexible,
  keywordstyle=\color{blue!60!black},commentstyle=\color{green!40!black}}

\newtheorem{theorem}{Theorem}
\newtheorem{lemma}{Lemma}
\newtheorem{definition}{Definition}
\newtheorem{invariant}{Invariant}

\title{\textbf{Lux DEX Cross-Chain Atomic Settlement}\\
\large Hash Time-Locked Contracts for Non-Custodial RFQ Settlement of Real BTC/ETH}
\author{Lux Cryptography}
\date{2026}

\begin{document}
\maketitle

\begin{abstract}
We replace the Lux DEX custodial lock-and-mint bridge and the \texttt{dexvm}
custodial atomic proxy with a genuine two-chain \emph{hash time-locked contract}
(HTLC) atomic swap. The settlement venue degenerates to a pure RFQ
\emph{coordinator}: it matches a user (taker) to a CEX-backed filler (maker),
relays a hashlock and on-chain lock proofs, and \emph{never holds keys, never
holds funds, and is never a recipient or refund party of any lock}. Settlement
either completes in full on both chains or both legs refund to their original
owners; no third party can mint, seize, or strand funds. We give the protocol,
the explicit timeout-ordering rule that guarantees atomicity, the on-Lux HTLC
precompile interface (selector ABI, state layout, events, per-function
verification), the Bitcoin~/~EVM counterparty-leg scripts, a threat model with
mitigations for every classical atomic-swap attack (free option, griefing,
preimage front-running, dust/fee griefing, reorg), and the migration that makes
every custodial mint path unreachable.
\end{abstract}

\section{Threat Model and Goals}

\paragraph{Adversary.} Any subset of: the venue/sequencer, the maker, a
network-level observer/front-runner, a chain reorg adversary bounded by $k$
blocks, and a malicious validator. The adversary is computationally bounded and
cannot break SHA-256 (preimage / collision resistance) or the signature schemes
of the two settlement chains (ECDSA/secp256k1 on Bitcoin and Lux~C-Chain).

\paragraph{Security goals.}
\begin{enumerate}[label=(G\arabic*),leftmargin=3.5em]
  \item \textbf{Non-custody.} No party other than the asset's owner can move it.
        The venue holds no key and appears in no lock as recipient or refund
        address.
  \item \textbf{Atomicity.} For an honest party that follows the protocol and is
        online within its claim window, the swap either completes on both legs
        or refunds on both legs. An honest party never ends with strictly less
        than it started (modulo agreed swap terms and gas).
  \item \textbf{No mint.} No wrapped asset is created; every unit paid out on
        either leg was locked by its owner on that same leg. There is no
        cross-chain supply.
  \item \textbf{Censorship/strand resistance.} Neither counterparty nor the
        venue can lock the other's funds indefinitely; every lock has an
        unconditional refund after its timeout.
\end{enumerate}

\section{Roles and the RFQ Flow}

\begin{definition}[Parties]
\textbf{$M$ (maker / filler)}: a CEX-backed market maker holding inventory on
both chains, who answers RFQ quotes and hedges on a centralized venue.
\textbf{$U$ (user / taker)}: holds the asset to be sold on one chain and wants
the asset on the other. \textbf{$V$ (venue)}: RFQ coordinator; matches, relays
messages and lock proofs; holds no funds, no keys, no preimage.
\end{definition}

We swap $M$'s asset on chain $A$ for $U$'s asset on chain $B$ (e.g. $A=$
Lux~C-Chain USDC, $B=$ Bitcoin). The assignment is deliberate (\S\ref{sec:to}):
\textbf{the maker locks first; the user generates the secret and is the
secret-holder.} This puts \emph{all} structural optionality on the professional
maker (who prices it into the quote) and gives the retail user the protected
position.

\begin{enumerate}[leftmargin=2.2em]
  \item \textbf{Quote.} $V$ publishes fill options (aggregated CEX prices /
        oracle). $U$ requests a quote for size $q$. $M$ returns a firm quote
        $(\text{rate}, \text{amount}_A, \text{amount}_B)$ valid for a short
        window $w$ (e.g. $w=30$\,s), signed by $M$.
  \item \textbf{Commit hashlock.} $U$ draws a uniform secret
        $s \xleftarrow{\$} \{0,1\}^{256}$, computes $h=\mathrm{SHA256}(s)$, and
        sends $h$ (with the accepted quote) to $M$ via $V$. $s$ never leaves
        $U$ until the claim.
  \item \textbf{Leg~A lock (maker first).} $M$ locks on chain $A$:
        $\mathrm{HTLC}_A(h,\ \text{recipient}=U,\ \text{refund}=M,\
        \text{asset}_A,\ \text{amount}_A,\ T_A)$. $M$ publishes the lock
        proof (tx id) to $V$; $V$ relays to $U$.
  \item \textbf{Verify + Leg~B lock (taker second).} $U$ waits for
        $K_A$ confirmations of $\mathrm{HTLC}_A$, checks $h$, amount, recipient
        $=U$, and $T_A$. If anything is wrong, $U$ does nothing and $M$ refunds
        at $T_A$. Else $U$ locks on chain $B$:
        $\mathrm{HTLC}_B(h,\ \text{recipient}=M,\ \text{refund}=U,\
        \text{asset}_B,\ \text{amount}_B,\ T_B)$ with $T_B > T_A + \Delta$
        (\S\ref{sec:to}). $U$ publishes the proof.
  \item \textbf{Claim (reveals $s$).} $U$ claims $\mathrm{HTLC}_A$ on chain $A$
        with $\mathrm{claim}(s)$, receiving $\text{amount}_A$. This publishes
        $s$ on chain $A$. $U$ must claim before $T_A$.
  \item \textbf{Counter-claim.} $M$ reads $s$ from chain $A$ (the
        \texttt{Claimed} event / witness) and claims $\mathrm{HTLC}_B$ on chain
        $B$ with $\mathrm{claim}(s)$ before $T_B$, receiving $\text{amount}_B$.
  \item \textbf{Refund (only off the happy path).} If $U$ never reveals $s$,
        $\mathrm{HTLC}_A$ refunds to $M$ at $T_A$ and $\mathrm{HTLC}_B$ refunds
        to $U$ at $T_B$. Both sides whole.
\end{enumerate}

After step~3 the venue is irrelevant to safety: steps 4--7 are peer-to-peer
on-chain operations that $U$, $M$, or their watchtowers execute directly. $V$
crashing or lying can only deny service or relay a stale quote --- never move
funds.

\section{The Timeout-Ordering Rule (atomicity invariant)}
\label{sec:to}

\begin{definition}[Secret-holder]
The secret-holder is the party that generates $s$ and whose claim first
publishes $s$. Here it is the user $U$.
\end{definition}

\begin{theorem}[Timeout ordering]
\label{thm:to}
Let the secret-holder claim the lock with the \emph{earlier} timeout and let the
counterparty claim the lock with the \emph{later} timeout. Atomicity (G2) holds
for online honest parties iff
\[
  T_{\text{late}} \;>\; T_{\text{early}} + \Delta,
  \qquad
  \Delta \;\ge\; R_{\text{late}} + C_{\text{late}},
\]
where $R_{\text{late}}$ is the reorg-safety depth (expressed in time) of the
chain holding the later-expiry leg and $C_{\text{late}}$ is the worst-case time
to confirm a claim there.
\end{theorem}

In our assignment the secret-holder $U$ claims $\mathrm{HTLC}_A$ (timeout $T_A$,
earlier) and the counterparty $M$ claims $\mathrm{HTLC}_B$ (timeout $T_B$,
later). Hence the normative rule:
\begin{center}\fbox{\parbox{0.92\linewidth}{\centering
\textbf{Maker's leg expires first; user's leg expires last:}\quad
$T_B > T_A + \Delta$.\\
The secret-holder's claim target ($\mathrm{HTLC}_A$) has the shorter timeout;
the counterparty's claim target ($\mathrm{HTLC}_B$) has the longer timeout.}}
\end{center}

\begin{proof}[Atomicity sketch]
Two outcomes only.
\emph{(i) $U$ reveals $s$ before $T_A$.} Then $s$ is public on chain $A$ at some
$t^\ast < T_A < T_B-\Delta$. $M$, online, claims $\mathrm{HTLC}_B$ in
$[t^\ast, T_B)$; the window has length $> \Delta \ge R_{\text{late}} +
C_{\text{late}}$, so even under a $k$-block reorg of chain $B$ the claim
confirms before $T_B$. Both legs claimed: complete.
\emph{(ii) $U$ does not reveal $s$ before $T_A$.} $\mathrm{HTLC}_A$ refunds to
$M$ at $T_A$. Since $s$ is unknown to $M$ (SHA-256 preimage resistance) and to
everyone else, $\mathrm{HTLC}_B$ cannot be claimed and refunds to $U$ at $T_B$.
Both legs refunded: rollback.
A reorg that removes $U$'s revealing claim on chain $A$ does \emph{not}
un-reveal $s$: once observed, $s$ is known, so $M$ can still claim
$\mathrm{HTLC}_B$. Recipient binding (\S\ref{sec:precompile}) ensures the
revealed $s$ enables payment only to the fixed recipients, never to a
front-runner.
\end{proof}

\paragraph{Concrete parameters.}
With $A=$ Lux (sub-second deterministic finality, $R_A\approx 0$) and $B=$
Bitcoin: set $K_A=$ Lux-finalized (1 block), $K_B=6$ confirmations,
$T_A = t_0 + 40\text{ min}$, $\Delta \ge 2\text{ h}$ (covering a 6-block BTC
reorg $\approx 60$\,min plus claim confirmation), $T_B = T_A + \Delta =
t_0 + 2\text{h}\,40\text{m}$. Bitcoin timeouts use \texttt{OP\_CHECKLOCKTIMEVERIFY}
on median-time-past; Lux timeouts use block timestamp. When both legs are EVM,
$\Delta$ shrinks to a few minutes.

\section{On-Lux HTLC Precompile (LP-90A0, \texttt{0x{\dots}90A0})}
\label{sec:precompile}

The HTLC is an \emph{orthogonal new primitive} with its own slot, not bolted
onto the AMM (\texttt{0x9010}) or CLOB (\texttt{0x9020}). It lives in a new
package \texttt{github.com/luxfi/precompile/swap} registered at
\texttt{0x0000{\dots}90A0} (\texttt{LXSwapAddress}, ``LP-90A0 LXSwap''), the next
free address after LP-9090 (LXSettlement). It reuses the proven-safe
\texttt{transferFrom}-delta custody idiom from \texttt{erc20\_vault.go} and the
standard \texttt{contract.StatefulPrecompiledContract} \texttt{Run} dispatch.

\subsection{Hash function: SHA-256 (cross-chain, not keccak)}
The hashlock MUST be $\mathrm{SHA256}$, matching Bitcoin's \texttt{OP\_SHA256}
and the EVM \texttt{0x02} precompile, so a single preimage $s$ unlocks both
legs. The preimage is fixed at exactly 32 bytes (Bitcoin witness compatibility).
\emph{The keccak commit--reveal in \texttt{hooks.go} is for the same-chain MEV
hook and is NOT reused here.}

\subsection{Selectors and Go interface}
\begin{lstlisting}[language=Go]
// Package swap is the Lux DEX cross-chain HTLC atomic-swap settlement
// precompile at LP-90A0. Non-custodial: funds move only to the recipient or
// refund address fixed at lock time; no owner, no pause, no upgrade.
package swap

const LXSwapAddress = "0x00000000000000000000000000000000000090A0"

// 4-byte selectors = BigEndian uint32 of keccak256(sig)[:4].
const (
    SelectorLock    uint32 = 0x........ // lock(bytes32,address,address,address,uint256,uint64)
    SelectorClaim   uint32 = 0x........ // claim(bytes32,bytes32)
    SelectorRefund  uint32 = 0x........ // refund(bytes32)
    SelectorGetSwap uint32 = 0x........ // getSwap(bytes32) view
    SelectorGetPreimage uint32 = 0x.... // getPreimage(bytes32) view  (hashlock -> s)
)

type SwapContract struct{} // holds NO mutable Go state; all state in StateDB.

func (c *SwapContract) Run(
    accessibleState contract.AccessibleState,
    caller  common.Address,
    addr    common.Address,
    input   []byte,
    suppliedGas uint64,
    readOnly bool,
) (ret []byte, remainingGas uint64, err error)
\end{lstlisting}

Internal handlers (mirroring the \texttt{0x9010} pattern):
\begin{lstlisting}[language=Go]
// lock(hashlock, recipient, refund, asset, amount, timeout) -> swapId
//   asset == common.Address{} (zero) means native LUX.
//   For native, requires the call-value seam (see Prerequisite below):
//   accessibleState.GetCallValue() == amount. For ERC20, performs a real
//   transferFrom and credits the OBSERVED balance delta (== amount), exactly
//   as erc20_vault.go safeTransferTokenFrom.
func (c *SwapContract) runLock(st contract.AccessibleState, caller common.Address,
    data []byte, gas uint64, readOnly bool) (swapId [32]byte, ret []byte, rem uint64, err error)

// claim(swapId, preimage) -> bool ; pays the STORED recipient (caller-agnostic).
func (c *SwapContract) runClaim(st contract.AccessibleState, caller common.Address,
    data []byte, gas uint64, readOnly bool) (ret []byte, rem uint64, err error)

// refund(swapId) -> bool ; pays the STORED refund address after timeout.
func (c *SwapContract) runRefund(st contract.AccessibleState, caller common.Address,
    data []byte, gas uint64, readOnly bool) (ret []byte, rem uint64, err error)
\end{lstlisting}

\subsection{State layout (StateDB on the precompile address)}
All persistent state lives in the EVM trie under \texttt{0x{\dots}90A0} using the
mapping-slot idiom $\textsf{slot}(id,f)=\mathrm{keccak256}(id \,\|\, f)$:
\begin{center}
\begin{tabular}{@{}lll@{}}
\toprule
Field & Slot key & Contents \\
\midrule
status      & $\mathrm{keccak}(\text{swapId}\,\|\,0)$ & 0 None / 1 Locked / 2 Claimed / 3 Refunded \\
hashlock    & $\mathrm{keccak}(\text{swapId}\,\|\,1)$ & $h=\mathrm{SHA256}(s)$ \\
recipient   & $\mathrm{keccak}(\text{swapId}\,\|\,2)$ & payout address on claim \\
refund      & $\mathrm{keccak}(\text{swapId}\,\|\,3)$ & payout address on refund \\
asset       & $\mathrm{keccak}(\text{swapId}\,\|\,4)$ & ERC20 addr, or $0$ for native \\
amount      & $\mathrm{keccak}(\text{swapId}\,\|\,5)$ & uint256 \\
timeout     & $\mathrm{keccak}(\text{swapId}\,\|\,6)$ & uint64 unix seconds \\
preimage    & $\mathrm{keccak}(\text{swapId}\,\|\,7)$ & $s$, zero until claimed \\
preimageIdx & $\mathrm{keccak}(\texttt{"pre"}\,\|\,h)$ & $s$ keyed by hashlock (cross-leg lookup) \\
reserve     & $\mathrm{keccak}(\texttt{"res"}\,\|\,\text{asset})$ & $\sum$ live locked amount for asset \\
nonce       & $\mathrm{keccak}(\texttt{"nonce"})$ & monotonic lock counter \\
\bottomrule
\end{tabular}
\end{center}
\[
  \text{swapId}=\mathrm{keccak256}\big(h \,\|\, \text{recipient} \,\|\,
  \text{refund} \,\|\, \text{asset} \,\|\, \text{amount} \,\|\, \text{timeout}
  \,\|\, \text{caller} \,\|\, \text{nonce}\big).
\]

\begin{invariant}[Reserve / conservation]
\label{inv:res}
For every asset, $\mathrm{balanceOf}(\texttt{0x{\dots}90A0}) \ge
\sum_{\text{Locked swaps of asset}} \text{amount} = \textsf{reserve}[\text{asset}]$.
Each \texttt{lock} increments \textsf{reserve} by the \emph{observed} inbound
delta; each \texttt{claim}/\texttt{refund} decrements it by exactly the swap's
amount and pays out the same amount. No code path mints.
\end{invariant}

\subsection{Per-function verification (normative)}
\textbf{lock:}
\begin{enumerate}[noitemsep,leftmargin=2em]
  \item \texttt{!readOnly}; else revert.
  \item $h\neq 0$; recipient $\neq 0$; refund $\neq 0$.
  \item $\text{amount} \ge \textsf{MinSwapAmount}$ (dust guard, per asset).
  \item $T_{0}+\textsf{MinTimeout} \le \text{timeout} \le T_{0}+\textsf{MaxTimeout}$,
        where $T_0=$ block timestamp. (Absolute sanity bounds; the relative
        $T_B>T_A+\Delta$ ordering is chosen by the parties off-chain and
        enforced by each party's own verification before locking.)
  \item Move funds in: native $\Rightarrow$ \texttt{GetCallValue()}$=$amount;
        ERC20 $\Rightarrow$ \texttt{transferFrom(caller,this,amount)} and require
        observed $\Delta\mathrm{balanceOf}=$ amount (fee-on-transfer safe).
  \item $\textsf{status}[\text{swapId}]=\text{None}$ (no collision); set fields;
        $\textsf{reserve}[\text{asset}]\mathrel{+}=$ amount; \texttt{nonce}{+}{+};
        emit \texttt{Locked}.
\end{enumerate}
\textbf{claim:}
\begin{enumerate}[noitemsep,leftmargin=2em]
  \item \texttt{!readOnly}; load swap; require $\textsf{status}=\text{Locked}$.
  \item require $T_0 < \text{timeout}$ (not expired).
  \item require $\mathrm{len}(\text{preimage})=32$ and
        $\mathrm{SHA256}(\text{preimage})=h$.
  \item pay the \emph{stored} recipient (caller-agnostic: a watchtower may
        submit); $\textsf{reserve}\mathrel{-}=$ amount; store preimage and
        $\textsf{preimageIdx}[h]$; $\textsf{status}=\text{Claimed}$; emit
        \texttt{Claimed}(swapId, $h$, preimage, recipient, amount).
\end{enumerate}
\textbf{refund:}
\begin{enumerate}[noitemsep,leftmargin=2em]
  \item \texttt{!readOnly}; load; require $\textsf{status}=\text{Locked}$.
  \item require $T_0 \ge \text{timeout}$.
  \item pay the \emph{stored} refund address; $\textsf{reserve}\mathrel{-}=$
        amount; $\textsf{status}=\text{Refunded}$; emit \texttt{Refunded}.
\end{enumerate}
The boundary is clean and disjoint: \texttt{claim} requires
$T_0<\text{timeout}$, \texttt{refund} requires $T_0\ge\text{timeout}$; a swap is
never both claimable and refundable, and (status$=$Locked once) never twice.

\paragraph{Events.}
\texttt{Locked(bytes32 indexed swapId, bytes32 indexed hashlock, address
recipient, address refund, address asset, uint256 amount, uint64 timeout)};
\texttt{Claimed(bytes32 indexed swapId, bytes32 indexed hashlock, bytes32
preimage, address recipient, uint256 amount)};
\texttt{Refunded(bytes32 indexed swapId, address refund, uint256 amount)}.
The \texttt{Claimed} event is the on-chain secret-relay: the counterparty's
watchtower subscribes by \texttt{hashlock} and reads \texttt{preimage}.

\paragraph{No admin surface.} \texttt{0x{\dots}90A0} has \emph{no} owner, pause,
freeze, upgrade, or set-balance selector. The only state transitions are
lock~$\to$~claim and lock~$\to$~refund, both permissionless and bounded by the
timeout. This is what makes (G1) and (G4) structural rather than policy.

\paragraph{Prerequisite (call-value seam).} Native-LUX locks require the
precompile to observe the EVM call value. The current
\texttt{StatefulPrecompiledContract.Run} signature carries no value (this is the
root of the \texttt{0x9010} native unbacked-mint finding). Add
\texttt{AccessibleState.GetCallValue() *big.Int} (or pass value into \texttt{Run})
in \texttt{contract/interfaces.go}, wired from the EVM call frame. Until that
lands, support ERC20/wrapped assets only (the \texttt{transferFrom}-delta path is
already value-safe).

\section{Counterparty-Leg Scripts}

\subsection{Bitcoin (P2WSH HTLC)}
\begin{lstlisting}
OP_IF
    OP_SHA256 <h> OP_EQUALVERIFY
    <recipient_pubkey> OP_CHECKSIG          ; claim path: witness = <sig> <s> 1
OP_ELSE
    <T_locktime> OP_CHECKLOCKTIMEVERIFY OP_DROP
    <refund_pubkey> OP_CHECKSIG             ; refund path: witness = <sig> 0
OP_ENDIF
\end{lstlisting}
$h=\mathrm{SHA256}(s)$, $|s|=32$. Timeout via absolute \texttt{OP\_CLTV}
(median-time-past). The refund spend sets \texttt{nLockTime}$\ge T$ and a
non-final sequence. Require $K_B=6$ confirmations before acting on a seen lock
and before treating a claim as final.

\subsection{EVM counterparty (Ethereum / another L2)}
Deploy the same HTLC logic as a Solidity contract (or this precompile on a
Lux-derived chain). Hashlock uses the \texttt{0x02} SHA-256 precompile, never
keccak, to remain preimage-compatible with the Bitcoin and Lux legs.

\section{Threat Model: Attacks and Mitigations}

\begin{center}\small
\begin{tabular}{@{}p{0.235\linewidth}p{0.715\linewidth}@{}}
\toprule
Attack & Mitigation \\
\midrule
\textbf{Free option / optionality.} Some party can wait and walk if price moves.
& Unavoidable in 2-of-2 HTLC; we \emph{allocate} it to the professional maker.
Maker locks first for a \emph{short} $T_A$ (e.g. 40\,min) and is the committed
party; the maker prices the option premium into the RFQ spread (it hedges on a
CEX). The user (taker) holds the option and the last-look --- the protected
side. Firm-quote window $w$ is seconds, bounding price drift. \\[2pt]
\textbf{Maker griefing / capital lockup.} A spam taker makes many makers lock,
then never proceeds, freezing maker capital until $T_A$.
& Bounded by the short $T_A$ (auto-refund). Venue-level rate-limiting and maker
reputation on takers. Optional hardening: taker posts the hashlock plus a small
slashable dust bond on the Lux leg \emph{before} the maker locks; the bond is
forfeit only if the taker fails to lock leg~B after the maker locks. The bond is
$\ll$ trade size, so it grants the maker no real option over the user. \\[2pt]
\textbf{Preimage front-running / observation.} An observer copies $s$ from the
mempool/witness when $U$ claims.
& Harmless: every HTLC binds payout to a \emph{fixed} recipient
(\texttt{CHECKSIG} to a specific pubkey on BTC; stored \texttt{recipient} on
Lux). Knowing $s$ lets one pay only the predetermined recipients, so a
front-runner at most completes the swap \emph{for} an honest party. Theft via
preimage requires an unbound ``anyone-with-$s$'' recipient, which we never use. \\[2pt]
\textbf{Fee / dust griefing.} A lock so small that the claim gas/fee exceeds the
output makes claiming irrational, stranding value.
& \textsf{MinSwapAmount} per asset on the Lux leg; Bitcoin outputs kept above
the dust limit and sized so $\text{amount}\gg$ claim fee. Refund is likewise
always economically rational. \\[2pt]
\textbf{Reorg on BTC/ETH leg.} A lock or claim is reorged out.
& Wait $K_A,K_B$ confirmations before treating a lock as real or a claim as
final; the safety buffer $\Delta \ge R_{\text{late}}+C_{\text{late}}$
(Thm.~\ref{thm:to}) covers a $k$-block reorg of the later-expiry leg. A reorg of
the secret-revealing claim does not un-reveal $s$. \\[2pt]
\textbf{Venue compromise / lying coordinator.} $V$ is malicious or offline.
& $V$ holds no key, no funds, and is no lock's recipient/refund. Worst case:
denial of service or a stale/bad quote (which the user can reject before
locking). Once both legs are locked, settlement is peer-to-peer; $V$ cannot move
funds. \\[2pt]
\textbf{Wrong-asset / amount / recipient in a lock.} Maker locks the wrong terms.
& The taker verifies $h$, asset, amount, recipient$=U$, and $T_A$ on-chain
\emph{before} locking leg~B (step~4). On mismatch the taker never locks; the
maker's lock simply refunds. \\[2pt]
\textbf{Hash-function mismatch across chains.} keccak vs SHA-256.
& Normative: SHA-256 hashlock on \emph{every} leg, $|s|=32$. The keccak
commit--reveal in \texttt{hooks.go} is a separate same-chain primitive and is
not used. \\
\bottomrule
\end{tabular}
\end{center}

\begin{theorem}[Non-custody]
Under the protocol, no party other than an asset's owner can transfer it, and
the venue can never transfer any asset.
\end{theorem}
\begin{proof}
Each unit sits in exactly one HTLC whose only exits are \texttt{claim} (to the
fixed recipient, gated by knowledge of $s$ with $\mathrm{SHA256}(s)=h$) and
\texttt{refund} (to the fixed refund address, gated by $T_0\ge T$). $V$ is
neither recipient nor refund of any HTLC and never learns $s$ before it is
public, so $V$ can call only \texttt{claim} with a public $s$ (paying the honest
recipient) or \texttt{refund} (paying the honest refund address) --- in both
cases value flows to the owner-designated address, not to $V$. The maker can
claim leg~B only after $U$ has already claimed leg~A (which revealed $s$),
i.e. only after $U$ received its side; by Invariant~\ref{inv:res} no path mints.
\end{proof}

\section{Migration: Removing Every Custodial Mint Path}

The audit found four reachable custodial/mint vectors. Each is removed or gated:
\begin{enumerate}[leftmargin=2em]
  \item \textbf{Bridge lock-and-mint (\texttt{precompile/bridge}).} Delete /
        disable \texttt{gateway.go} \texttt{CompleteBridge}, \texttt{InitiateBridge}
        and the liquidity pools (the mint-on-no-op-signature path). Remove the
        no-op verifiers \texttt{gateway.go:verifySignature} and
        \texttt{signer.go:VerifyThresholdSignature}. The MPC signer set is
        retained, if at all, only for non-custodial coordination (it can sign
        \emph{advisory} RFQ attestations but holds no mint authority).
  \item \textbf{\texttt{dexvm} custodial proxy (\texttt{luxfi/chains/dexvm}).}
        Reject client-submitted \texttt{TxImport}/\texttt{TxExport} at admission
        (\texttt{chainvm.go:SubmitTx}~/~\texttt{feegate.go:gateUserTxBytes}) and
        in \texttt{vm.go:processTx}; the atomic shared-memory export leg
        (\texttt{atomic.go:executeExport}) is no longer a settlement path.
        Trades settle via the HTLC precompile between user and filler directly.
  \item \textbf{EVM custody native-mint (\texttt{precompile/dex/pool\_manager.go}).}
        Retire the custodial \texttt{deposit}/\texttt{withdraw} ZAP-credit rail
        (the \texttt{0x9010} native-balance unbacked-mint at
        \texttt{pool\_manager.go:1396}). The CLOB/RFQ venue keeps matching; final
        asset settlement moves to \texttt{0x{\dots}90A0}.
  \item \textbf{Stubbed light clients (\texttt{node/vms/chainadapter}).} No
        longer on any settlement path: the HTLC legs are verified by each party's
        own full node, not by the in-process stubs. Remove from the mint path;
        if retained for watchtower convenience they must do real SPV before any
        reliance.
\end{enumerate}
After migration the only cross-chain value movement is the HTLC pair; grepping
for a reachable mint (\texttt{AddBalance} not backed by an observed inbound
transfer, or a shared-memory \texttt{PutRequest} from client input) returns
nothing in the settlement path.

\section{Conclusion}
The construction makes non-custody (G1) and strand-resistance (G4) structural
properties of the on-chain HTLCs rather than policies enforced by a trusted
operator, and reduces atomicity (G2) and no-mint (G3) to SHA-256 preimage
resistance plus the timeout-ordering rule $T_B > T_A + \Delta$. The venue is
demoted to a coordinator that can delay or misquote but never steal, freeze, or
mint. This is the minimal trusted-computing base for real-asset DEX settlement.

\end{document}
